\section{Examples}\label{sec:examples}

We illustrate the divisor construction in rank two and the global projective line
bundle construction for a proper root ideal in rank three.

\subsection{The divisor construction in rank two}

Let $X=\SX_{(0,1)}$ be the compactification of $T^*\PP^1$, and set $\Lambda_{G\times S}=\Z[c_1,c_2,\epsilon]$. By Corollary~\ref{cor:equiv_coho_xdelta_G}, the equivariant cohomology ring is presented as
\[
 H^*_{G\times S}(X;\Z)=\Lambda_{G\times S}[x_1,x_2]/(F_1,F_2),
\]
where
\[
 F_1=x_1^2+c_1x_1+c_2,\qquad
 F_2=x_2^2+c_1x_2+c_2-\epsilon(x_1+x_2+c_1),
\]
with the $\Lambda_{G\times S}$-basis $\{1,x_1,x_2,x_1x_2\}$.
Set $\ell=x_1+x_2+c_1$ and $E=x_2-x_1-\epsilon$; then
$F_2-F_1=E\ell$.
For the smooth divisor $j:Y=\SX_{(0,0)}\cong\PP^1\hookrightarrow X$, we have
\[
 H^*_{G\times S}(Y;\Z)=\Lambda_{G\times S}[x_1,x_2]/(F_1,\ell),\qquad
 j_*[g]=[Eg].
\]
Notice that the divided difference operator $\partial_1$ does not descend to $H^*_{G\times S}(X;\Z)$, because
$\partial_1(F_1)=\ell=x_1+x_2+c_1$ is nonzero in $H^*_{G\times S}(X;\Z)$.
After restriction to the divisor $Y$, however, the relation $\ell=0$ holds, so $\partial_1(F_1)\in(F_1,\ell)$ and the push--pull operator $\pi^*\pi_*$ along
$\pi:Y\to\pt$ acts on $H^*_{G\times S}(Y;\Z)$ as $\partial_1$.
Composing with the Gysin pushforward $j_*$ and adding the identity therefore gives, for every $g\in \Lambda_{G\times S}[x_1,x_2]$,
\[
 (\id+j_*\pi^*\pi_*j^*)[g]=[(\id+E\partial_1)g]=[\BS_1g]\qquad\text{in }H^*_{G\times S}(X;\Z).
\]

\subsection{A projective line bundle for a proper root ideal in rank three}

Let $n=3$ and $\rfunc=(0,1,1)$, so that
$\Psi=\{\epsilon_1-\epsilon_2,\epsilon_1-\epsilon_3\}$.
The simple reflection $s_2$ preserves $\Psi\cup(-\Psi)$, while $\alpha_2=\epsilon_2-\epsilon_3\notin\Psi$.
Set $X=\SX_{(0,1,1)}$, $B_2=\SX_{(0,1,1),2}^P$, $\Lambda_{G\times S}=\Z[c_1,c_2,c_3,\epsilon]$, $f(z)=z^3+c_1z^2+c_2z+c_3$, and
\[
 g(u,v)=\frac{f(u)-f(v)}{u-v}
       =u^2+uv+v^2+c_1(u+v)+c_2.
\]
Then $H^*_{G\times S}(X;\Z)\cong\Lambda_{G\times S}[x_1,x_2,x_3]/(D_1,D_2,D_3)$, where
\[
 \begin{aligned}
 D_1&=f(x_1),\\
 D_2&=f(x_2)-\epsilon g(x_1,x_2),\\
 D_3&=g(x_2,x_3)-\epsilon(x_1+x_2+x_3+c_1).
 \end{aligned}
\]
Because the condition $tL_3\subset L_1$ imposes no constraint on $L_2$, forgetting $L_2$
defines a global $\PP^1$-bundle $\pi_2:X\to B_2$.
The relation ideal $(D_1,D_2,D_3)$ is stable under $\partial_2$, since
\[
 \partial_2D_1=0,\qquad \partial_2D_2=D_3,\qquad \partial_2D_3=0,
\]
and stability follows from the twisted Leibniz rule.
The push--pull operator $\pi_2^*\pi_{2*}$ acts on $H^*_{G\times S}(X;\Z)$ by $\partial_2$, and the twisted relative cotangent line bundle has Euler class
$e_{G\times S}(\Omega^1_{X/B_2}\otimes\C_{-1})=x_3-x_2-\epsilon$.
Thus the second construction of Theorem~\ref{thm:geometric_hecke} gives
\[
 \id+e_{G\times S}(\Omega^1_{X/B_2}\otimes\C_{-1})\cup\pi_2^*\pi_{2*}
 =\id+(x_3-x_2-\epsilon)\partial_2=\BS_2.
\]
